\documentclass[aspectratio=169,10pt]{beamer}
\usetheme{Madrid}
\usecolortheme{default}

\usepackage{hyperref}
\usepackage{tikz}
\usepackage{amsmath,amssymb}

% --- Custom Colors ---
\definecolor{unibblue}{RGB}{0,51,102}
\definecolor{unibgold}{RGB}{212,175,55}
\definecolor{unibgreen}{RGB}{0,128,0}

\setbeamercolor{palette primary}{bg=unibblue,fg=white}
\setbeamercolor{palette secondary}{bg=unibblue!85,fg=white}
\setbeamercolor{palette tertiary}{bg=unibblue!70,fg=white}
\setbeamercolor{structure}{fg=unibblue}
\setbeamercolor{block title}{bg=unibblue!85,fg=white}
\setbeamercolor{block body}{bg=unibblue!10,fg=black}
\setbeamercolor{alertblock title}{bg=unibgold!90,fg=black}
\setbeamercolor{alertblock body}{bg=unibgold!15,fg=black}

\title{Persamaan Diferensial}
\subtitle{Minggu 9: Pengantar Persamaan Diferensial Parsial (PDP) \& Analisis Deret Fourier}
\author{Novalio Daratha \& Muhammad Arfan}
\institute{Program Studi Teknik Elektro\\Universitas Bengkulu}
\date[\today]{2026 \\ \vspace{2mm}\small\textit{Tanggal Kompilasi: \today}}

\begin{document}

\begin{frame}
  \titlepage
\end{frame}

\begin{frame}{Capaian Pembelajaran (CPMK 4)}
  \begin{itemize}
    \item \textbf{CPMK 4.1:} Memahami perbedaan mendasar antara PDB (1 variabel bebas) dan PDP (multi-variabel bebas seperti ruang $x,y$ dan waktu $t$).
    \item \textbf{CPMK 4.2:} Memahami teorema ortogonalitas fungsi trigonometri.
    \item \textbf{CPMK 4.3:} Mampu menghitung koefisien Deret Fourier ($a_0, a_n, b_n$) untuk sinyal periodik non-sinusoidal.
    \item \textbf{CPMK 4.4:} Memanfaatkan simetri fungsi genap (*cosine series*) dan fungsi ganjil (*sine series*) untuk mempermudah integrasi.
  \end{itemize}
\end{frame}

\begin{frame}{Daftar Isi}
  \tableofcontents
\end{frame}

\section{Pengantar Persamaan Diferensial Parsial (PDP)}
\begin{frame}{Mengapa Kita Membutuhkan PDP?}
  Dalam fenomena fisik elektro, besaran seringkali bergantung pada \textbf{posisi spasial ($x,y,z$)} DAN \textbf{waktu ($t$)}.
  \begin{block}{Tiga Persamaan Klasik PDP dalam Teknik Elektro}
    \begin{enumerate}
      \item \textbf{Persamaan Gelombang 1D (Saluran Transmisi):}
      $$ \frac{\partial^2 u}{\partial t^2} = v^2 \frac{\partial^2 u}{\partial x^2} $$
      \pause
      \item \textbf{Persamaan Panas/Difusi 1D (Pemanasan Kabel):}
      $$ \frac{\partial u}{\partial t} = \alpha \frac{\partial^2 u}{\partial x^2} $$
      \pause
      \item \textbf{Persamaan Laplace 2D (Medan Elektrostatik):}
      $$ \frac{\partial^2 V}{\partial x^2} + \frac{\partial^2 V}{\partial y^2} = 0 $$
    \end{enumerate}
  \end{block}
\end{frame}

\section{Fondasi Deret Fourier Trigonometrik}
\begin{frame}{Prinsip Representasi Deret Fourier}
  \begin{alertblock}{Teorema Fourier}
    Setiap fungsi periodik $f(t)$ dengan periode $T = 2L$ yang memenuhi syarat Dirichlet dapat diuraikan menjadi jumlahan tak hingga gelombang harmonik sinus dan kosinus:
    $$ {\color{unibblue} f(t) = a_0 + \sum_{n=1}^{\infty} \left( a_n \cos\left(\frac{n\pi t}{L}\right) + b_n \sin\left(\frac{n\pi t}{L}\right) \right)} $$
  \end{alertblock}
  \pause
  \begin{block}{Rumus Euler-Fourier untuk Koefisien:}
    \begin{itemize}
      \item Nilai Rata-rata: $a_0 = \frac{1}{2L} \int_{-L}^{L} f(t) \, dt$
      \item Koefisien Kosinus: $a_n = \frac{1}{L} \int_{-L}^{L} f(t) \cos\left(\frac{n\pi t}{L}\right) \, dt$
      \item Koefisien Sinus: $b_n = \frac{1}{L} \int_{-L}^{L} f(t) \sin\left(\frac{n\pi t}{L}\right) \, dt$
    \end{itemize}
  \end{block}
\end{frame}

\section{Simetri Fungsi: Genap \& Ganjil}
\begin{frame}{Memanfaatkan Sifat Simetri Gelombang}
  \begin{columns}
    \begin{column}{0.5\textwidth}
      \begin{block}{Fungsi Genap: $f(-t) = f(t)$}
        \begin{itemize}
          \item Simetri terhadap sumbu-$y$.
          \item Hanya memuat suku kosinus!
          \item $b_n = 0$
          \item $a_n = \frac{2}{L}\int_{0}^{L} f(t)\cos\left(\frac{n\pi t}{L}\right) dt$
        \end{itemize}
      \end{block}
      \textit{Contoh: Gelombang segitiga genap, $\cos(\omega t)$.}
    \end{column}
    \begin{column}{0.5\textwidth}
      \pause
      \begin{block}{Fungsi Ganjil: $f(-t) = -f(t)$}
        \begin{itemize}
          \item Simetri terhadap titik asal $(0,0)$.
          \item Hanya memuat suku sinus!
          \item $a_0 = 0, \quad a_n = 0$
          \item $b_n = \frac{2}{L}\int_{0}^{L} f(t)\sin\left(\frac{n\pi t}{L}\right) dt$
        \end{itemize}
      \end{block}
      \textit{Contoh: Gelombang kotak simetris, gelombang gigi gergaji, $\sin(\omega t)$.}
    \end{column}
  \end{columns}
\end{frame}

\begin{frame}{Contoh Nyata: Penguraian Harmonik Gelombang Kotak}
  Tinjau gelombang kotak ganjil: $f(t) = -V_0$ untuk $-\pi < t < 0$, dan $f(t) = +V_0$ untuk $0 < t < \pi$ ($L = \pi$).
  \pause
  Karena fungsi ganjil: $a_0 = 0, a_n = 0$.
  \pause
  $$ b_n = \frac{2}{\pi} \int_{0}^{\pi} V_0 \sin(nt) \, dt = \frac{2V_0}{\pi} \left[ -\frac{\cos(nt)}{n} \right]_{0}^{\pi} = \frac{2V_0}{n\pi} (1 - (-1)^n) $$
  \pause
  \begin{itemize}
    \item Jika $n$ genap ($n=2,4,6,\dots$): $b_n = 0$.
    \item Jika $n$ ganjil ($n=1,3,5,\dots$): $b_n = \frac{4V_0}{n\pi}$.
  \end{itemize}
  \pause
  \begin{alertblock}{Deret Fourier Gelombang Kotak}
    $$ {\color{unibblue} f(t) = \frac{4V_0}{\pi} \left( \sin(t) + \frac{1}{3}\sin(3t) + \frac{1}{5}\sin(5t) + \dots \right)} $$
  \end{alertblock}
\end{frame}

\section{Kuis \& Rangkuman}
\begin{frame}{Kuis Interaktif 9}
  \textbf{Pertanyaan:} Sebuah inverter daya menghasilkan tegangan bolak-balik $v(t) = v(-t)$. Manakah koefisien Fourier yang bernilai pasti nol tanpa perlu dihitung?
  \pause
  \begin{block}{Jawaban:}
    Karena $v(t)$ adalah fungsi \textbf{genap} ($v(-t) = v(t)$), maka seluruh koefisien sinus:
    $$ {\color{unibblue} b_n = 0 \quad \text{untuk seluruh } n \ge 1} $$
    Sinyal tersebut murni dapat dinyatakan hanya dengan deret kosinus Fourier!
  \end{block}
\end{frame}

\begin{frame}{Rangkuman Materi Minggu 9}
  \begin{enumerate}
    \item PDP melibatkan turunan terhadap ruang dan waktu, membutuhkan ekspansi deret untuk memecahkan batas-batasnya.
    \item Deret Fourier membedah sinyal non-sinusoidal menjadi spektrum frekuensi harmonik.
    \item Pengenalan simetri genap/ganjil memangkas separuh proses kalkulasi integral secara efisien.
  \end{enumerate}
\end{frame}

\end{document}
